% !TEX root = ../../main.tex
\section{Tính cấp thiết của đề tài}
\label{sec:tinh_cap_thiet}
Trong kỷ nguyên số và bùng nổ thông tin hiện nay, số lượng ấn phẩm khoa học trên toàn cầu tăng trưởng với tốc độ chưa từng có. Mỗi tuần có hàng nghìn bài báo mới xuất hiện trên các cơ sở dữ liệu học thuật quốc tế uy tín như arXiv, Semantic Scholar, IEEE Xplore hay ACM Digital Library. Việc tiếp cận, đọc hiểu, phân loại và tổng hợp tri thức từ kho tài liệu khổng lồ này đặt ra áp lực quá tải thông tin rất lớn đối với các nhà nghiên cứu và sinh viên đại học.

Mặc dù sự xuất hiện của các mô hình ngôn ngữ lớn (LLM) mở ra tiềm năng to lớn trong việc hỗ trợ học thuật, việc ứng dụng các công cụ AI phổ quát hiện nay vào quy trình nghiên cứu khoa học chuyên sâu đang gặp phải bốn rào cản nghiêm trọng:
\begin{enumerate}[label=\textbf{(\arabic*)}]
    \item \textbf{Ảo giác trích dẫn và bịa đặt học thuật} (\textit{Citation Hallucination}): Các mô hình AI tạo sinh thông thường có xu hướng tự sáng tác ra tên tác giả, số định danh DOI, tên tạp chí và các số liệu thực nghiệm không có thật để trả lời câu hỏi của người dùng. Trong nghiên cứu khoa học, việc trích dẫn một nguồn tài liệu không tồn tại là vi phạm nghiêm trọng chuẩn mực liêm chính học thuật và đạo đức nghiên cứu \cite{lewis2020}.
    \item \textbf{Thiếu cơ chế kiểm chứng tọa độ trực quan trên tài liệu gốc:} Phần lớn các công cụ hỗ trợ đọc tài liệu hiện nay (như ChatPDF thông thường) sử dụng phương pháp bóc tách văn bản thô, làm mất hoàn toàn cấu trúc hai cột đặc trưng của các bài báo khoa học, loại bỏ các công thức toán học và chỉ trích dẫn số trang chung chung. Người nghiên cứu không thể xác minh được khẳng định của AI nằm chính xác tại dòng chữ nào, đoạn văn nào trên trang PDF gốc.
    \item \textbf{Sự phân mảnh của hệ sinh thái công cụ hỗ trợ nghiên cứu:} Quy trình nghiên cứu hiện nay bị chia cắt giữa nhiều phần mềm độc lập: học giả phải dùng Zotero / Mendeley để quản lý danh mục tài liệu, dùng Connected Papers / Litmaps để khám phá mạng lưới bài báo, dùng ChatPDF để tra cứu nhanh, và chuyển sang Overleaf để soạn thảo bản thảo \LaTeX. Sự thiếu liên kết này làm đứt gãy mạch tư duy và tiêu tốn đáng kể thời gian chuyển đổi ngữ cảnh của nhà nghiên cứu.
    \item \textbf{Rào cản kỹ thuật phức tạp khi biên soạn tài liệu khoa học bằng \LaTeX:} Để công bố công trình trên các hội nghị và tạp chí quốc tế hàng đầu, bản thảo bắt buộc phải tuân thủ nghiêm ngặt các khuôn mẫu (template) định dạng của IEEE Tran, ACM SIG hay Springer. Khi xảy ra lỗi biên dịch cú pháp trong \LaTeX, học giả thường mất hàng giờ để đọc các tệp nhật ký (\texttt{.log}) dài hàng nghìn dòng để dò tìm lỗi thủ công, làm giảm đáng kể hiệu suất hoàn thiện bài báo.
\end{enumerate}

Xuất phát từ các yêu cầu thực tiễn cấp bách đó, việc nghiên cứu và xây dựng một nền tảng trợ lý AI chuyên biệt, có khả năng đối chiếu bằng chứng tọa độ trực quan trên bản PDF gốc, xây dựng đồ thị tri thức 4 thực thể và tự động sửa lỗi biên dịch bản thảo \LaTeX\ là một đòi hỏi vô cùng cấp thiết nhằm nâng cao năng suất và chất lượng công bố khoa học quốc tế, đáp ứng đầy đủ quy chế đào tạo đại học và định hướng ứng dụng công nghệ thông tin tiên tiến \cite{bgd2021, humg2025}.
